\chapter{Introduction and Project Scope}
\label{ch:introduction}

\section{Background and Motivation}
\label{sec:intro-background}

Corrosion of steel reinforcement is a principal durability concern for
reinforced cementitious structural elements, and ferrocement — a thin,
densely mesh-reinforced cementitious composite — is no exception. Assessing
that corrosion typically requires either destructive sampling or
specialised instrumentation, both of which are costly to apply repeatedly
over an element's service life. Repeated photography of an exposed
surface is, by contrast, comparatively inexpensive and non-destructive,
and it raises a natural question: how much can be inferred about a
specimen's condition from its accumulated photographic record alone?

Answering that question requires distinguishing two quantities that are
easy to conflate but evidentially very different. Visible surface
corrosion is directly observable in a photograph and can, in principle, be
labelled and re-labelled from the same image at any time. Internal
structural condition, by contrast, is not observed longitudinally in the
photographic record at all: structural supervision is available only
through terminal measurements — wire-area loss and ultimate load —
obtained once per specimen, destructively, at the end of its exposure
period. This report's central technical thread, developed across the
chapters that follow, is this distinction: what a dense,
repeatedly-observed surface quantity can support scientifically is not
automatically what sparse, terminal structural measurements can support,
and the two must be evaluated, and reported, differently.

The physical specimens, their fabrication, their accelerated
chloride-exposure programme, and their terminal mechanical testing are
predecessor experimental work, external to the present activity,
documented in an earlier MSc thesis \citep{hossain_thesis} and a related
conference paper \citep{artiste2025}. The present activity builds on, and
uses, the experimental dataset generated by that preceding work; it does
not repeat the physical experiment. Some of the later image-feature work
in this activity is best described as a Python reconstruction and
extension of concepts from that predecessor workflow, rather than an
exact reproduction of it (Chapter~\ref{ch:activities}). Facts about the physical experiment are
attributed to that predecessor source throughout this report; facts about
how the resulting data has been organised, audited, modelled, and
interpreted are attributed to the present activity.

\section{Project Context and Data Basis}
\label{sec:intro-context}

The dataset comprises 48 ferrocement specimens tested across two
experimental campaigns, each photographed repeatedly over its exposure
period. Visible-corrosion labels are available throughout this
photographic record — in the current audited basis, 791 aligned, readable
observations out of 792 raw image files (Chapter~\ref{ch:dataset}) — while
terminal structural measurements exist for each specimen exactly once, at
the end of its exposure period. This asymmetry between dense, repeatedly
observed visible-corrosion supervision and sparse, single-timepoint
structural supervision is the single most consequential structural fact
about this dataset, and it is why visible-corrosion estimation and
structural-condition inference are treated as scientifically distinct
problems throughout this report rather than as two readings of the same
underlying task. Chapter~\ref{ch:dataset} presents the full dataset audit;
it is not repeated here.

\section{Project Objectives and Research Questions}
\label{sec:intro-objectives}

The activity was organised around five research questions, each answered,
to varying degrees of confidence, by the work reported in
Chapters~\ref{ch:activities} and~\ref{ch:results}.

\textbf{RQ1.} To what extent can visible surface corrosion be estimated
from repeated images, using classical, deep-embedding, and interpretable
representations in turn?

\textbf{RQ2.} Can hidden structural condition, or terminal structural
response specifically, be inferred reliably from the available
image-derived and metadata inputs?

\textbf{RQ3.} Do image-derived features provide incremental predictive
value beyond known specimen and exposure metadata for terminal ultimate
load?

\textbf{RQ4.} How robust are the resulting conclusions when evaluation
moves from a grouped specimen holdout to treatment-level and campaign-level
transfer tests?

\textbf{RQ5.} What can degradation-curve fitting and threshold-based
proxy-RUL screening legitimately contribute, given the absence of repeated
structural measurements and of any observed true lifetime outcome?

A later, separate workstream — preparing a leakage-safe dataset and
augmentation pipeline for four-class visible-corrosion severity
classification — is deliberately not framed as a sixth research question
here, since it addresses data preparation rather than a modelling
question, and no classifier has yet been trained on it
(Section~\ref{sec:intro-scope}).

\section{Scope of the Technical Activity}
\label{sec:intro-scope}

The activity reported here comprises: an audit and alignment of the raw
photographic and workbook data; a classical visible-corrosion baseline; a
frozen deep-image-embedding experiment; an engineered, interpretable image
feature representation; structural-condition feasibility modelling;
robustness evaluation under grouped, leave-one-treatment-out, and
leave-one-campaign-out regimes; degradation-curve fitting and proxy-RUL
screening; a refocused study of terminal ultimate load; preparation of an
augmented, specimen-grouped dataset for four-class classification; and an
audit of the project's own software and reproducibility practices.
Chapter~\ref{ch:activities} describes each of these in turn;
Chapter~\ref{ch:engineering} covers the last.

Several boundaries on what this report claims are stated here explicitly,
because they govern how every later chapter should be read. This report
does not establish, and does not claim, a validated remaining-useful-life
model: no true RUL supervision exists in the dataset, and the
proxy-RUL construct introduced in Chapter~\ref{ch:framework} is an
explicitly qualified screening quantity, not a lifetime prediction. It
does not claim that visible surface corrosion causally determines
structural capacity; where a statistical association between the two is
reported, it is reported together with the confounding structure that
limits how it may be interpreted. It does not produce an externally
validated or deployment-ready model of any kind. It does not include a
trained four-class classifier: that branch's data preparation is complete,
but no classifier has been trained, compared, or evaluated
(Chapter~\ref{ch:activities}). And it does not treat model-estimated
structural states at non-terminal observation weeks as measurements — they
are, throughout, the output of a model, not a repeated physical
observation.

\section{Main Contributions of the Activity}
\label{sec:intro-contributions}

The activity's principal, completed contributions are as follows.

\textbf{(A)} Audited and aligned a specimen-level analysis basis from the
raw photographic and workbook data, resolving several version-specific
discrepancies across the project's successive implementation phases
(Chapter~\ref{ch:dataset}).

\textbf{(B)} Established visible-corrosion estimation as the most
consistently learnable image-based task in this dataset, while
identifying specific label-adjacency and near-tautology risks that limit
how some individual benchmarks should be read
(Chapter~\ref{ch:activities}, Chapter~\ref{ch:results}).

\textbf{(C)} Quantified the substantially weaker and less stable
transferability of hidden structural prediction, and disentangled the
contribution of image-derived inputs from metadata/campaign-linked
predictive structure through explicit ablations and transfer tests
(Chapter~\ref{ch:activities}, Chapter~\ref{ch:results},
Chapter~\ref{ch:interpretation}).

\textbf{(D)} Refocused the structural research question onto terminal
ultimate load — one of the two directly measured terminal structural
endpoints, and the one selected for this analysis — and explicitly
tested, rather than assumed, whether images add predictive value beyond
metadata for it (Chapter~\ref{ch:activities}).

\textbf{(E)} Built and critically evaluated a degradation-curve and
proxy-RUL screening chain across two independently configured
generations, establishing why its contribution is methodological and
diagnostic rather than validated prognosis
(Chapter~\ref{ch:activities}, Chapter~\ref{ch:interpretation}).

\textbf{(F)} Prepared a leakage-safe, specimen-grouped, reproducible
dataset and augmentation pipeline for a four-class visible-corrosion
severity classification task, ready for subsequent model training
(Chapter~\ref{ch:activities}).

\textbf{(G)} Documented concrete reproducibility and software-engineering
issues in the project's own codebase, and maintained a traceable,
source-cited evidence record for every quantitative claim in this report
(Chapter~\ref{ch:engineering}).

None of these is presented as a novel publication-level scientific
contribution in its own right; each is a specific, bounded piece of work
completed during this activity, reported at the level of confidence its
own evidence supports.

Read together, the activity's findings are, at a high level: visible
surface corrosion is the most consistently learnable image-based quantity
available in this dataset; hidden structural inference is substantially
weaker, and strongly affected by sparse supervision and campaign-linked
design structure; images do not show a reliable pooled incremental
benefit over metadata for terminal ultimate load under the settings
tested here; and the proxy-RUL construct remains a useful exploratory
screening tool rather than a validated prognosis of remaining service
life. The evidence and qualifications behind each of these statements are
given in full in Chapters~\ref{ch:activities}--\ref{ch:interpretation}.

\section{Report Organisation}
\label{sec:intro-organisation}

Chapter~\ref{ch:dataset} audits the dataset and establishes the
supervision asymmetry this introduction has previewed.
Chapter~\ref{ch:framework} sets out the methodological framework —
evaluation regimes, representation strategies, and the
observed/model-derived/derived-screening distinction — shared across later
chapters. Chapter~\ref{ch:activities} reconstructs the project's
successive implementation phases chronologically.
Chapter~\ref{ch:results} integrates that evidence by scientific question
rather than by phase, and Chapter~\ref{ch:interpretation} explains what
the integrated evidence means and why it justified the research decisions
taken. Chapter~\ref{ch:engineering} documents the project's software and
reproducibility practices and their current limitations,
Chapter~\ref{ch:limitations} states the limitations of the evidence itself
by category, and Chapter~\ref{ch:status} gives the current project status
and the work this report recommends next.
